\section{Primitive and inferential details for controller qualification}\label{supp:displacement75}
For $x,y\in[0,1]$, the original transition in its first coordinate is
\[
 f_1(x,y,a,z)=\frac1{16}+\frac9{16}x+\frac18y-\frac1{16}xy+\frac12a+z;
\]
the second interchanges $x,y$, replaces $a/2$ by $a/4$, and replaces $z$ by $-z$. Every state-Jacobian entry is nonnegative. A column sum is at most $9/16+1/8=11/16$, and the action-column sum is $3/4$. Integration of the Jacobian along the line segment between two points gives the claimed $\ell^1$ bounds. With $0\le a\le1/4$ and $|z|\le1/32$, both coordinates remain in $[0,1]$. The same estimate therefore applies at every date.

The running state loss is
\[
 (x-5/8)^2+(y-5/8)^2+\tfrac14(x-y)^2+2(1/2-x-y)_+^2.
\]
Each component of its gradient has absolute value at most $5/4+1/2+2=15/4<4$. The derivative of $a^2+4a^4$ is at most $2(1/4)+16(1/4)^3=3/4\le13/16$ on the entire feasible action range. The terminal state loss doubles the first two squared deviations, giving the componentwise derivative bound $5/2+1/2+2=5$. A componentwise gradient bound is the Lipschitz constant for the $\ell^1$ distance. These estimates remain valid across the shortage boundary because the squared positive part is continuously differentiable.

The state part of the running cost is convex on the square. Its maximum is therefore bounded by its maximum at the four vertices, which is $82/64$ at $(0,0)$. The adjustment term is at most $5/64$, giving the running-cost bound $87/64$. The same convexity and vertex calculation bounds the terminal cost by $33/16$. Nonnegativity and $\beta=15/16$ thus give
\[
 0\le G^\pi\le H:=\frac{87}{64}(1+\beta)+\frac{33}{16}\beta^2
\]
for every complete feasible two-date controller. Taking $\min\{H,K_{\pi\nu}\}$ is valid for an absolute difference. None of these constants is estimated from the sampled cost rows.

For rectangular cells $C=\prod_j[l_j,u_j]$, $D=\prod_j[v_j,w_j]$, the closure distance is
\[
 \operatorname{dist}(C,D)=\sum_j\max\{0,l_j-w_j,v_j-u_j\}.
\]
The production support calculation uses exact integer-scaled versions of these distances. An independently written rational loop reconstructs every eligible pair and every datewise maximum. The controller at an exact acquisition boundary takes the upper cell, clipped at one. Using closed cells when constructing the support only enlarges the set of possible pairs and remains valid under this deterministic convention.

\subsection{Endpoint coverage and rounding}
For one estimand, assign failure probability $\gamma/(2M)$ to each endpoint statement. The empirical-Bernstein theorem bounds the mean of the lower endpoint from below and that of the upper endpoint from above, using the same unbiased sample variance as in \eqref{eq:cost-interval75}. Since $\mathbb E L\le\mu\le\mathbb E U$, a union bound over the $2M$ statements gives simultaneous coverage. This argument requires no independence across estimands. For the thirteen-estimand experiment $\gamma=1/100$ and $\ell=9$ suffice: an exact truncated exponential series exceeds $5200=4M/\gamma$, proving $\log(5200)<9$ without a floating logarithm.

The independent path evaluator uses arbitrary-precision integer endpoints at scale $2^{80}$. Every multiplication is divided with downward and upward integer rounding. It rounds its final endpoints outward to scale $2^{60}$, computes endpoint sums and squared sums as integers, and obtains exact rational means and variances. A square-root upper bound is obtained by integer square root followed, when necessary, by one upward unit on the sixty-bit grid. These roundings expand the displayed confidence interval rather than reduce it. The endpoint rounding is already inside the recorded mean enclosure; it must not be subtracted as if it were a statistical bias correction.

The production evaluator and independent evaluator share the frozen economic inputs, not numerical kernels. Production uses the factorized bilinear expression; the independent route uses its expanded polynomial. Both propagate whole boxes and include all intersected actor cells. Exact rational point and face tests are additional checks of the algebra. Nonoverlap of two claimed enclosures is a diagnostic failure. Overlap alone is not a proof that either enclosure contains the true path; validity also uses the stated primitive identities and outward arithmetic.

\subsection{Prices and complete work}
Multiplying a covered gain interval by the nonnegative quantity $D\lambda$ and subtracting the deterministic priced work preserves coverage, proving \eqref{eq:adoption75}. On the single simultaneous event this implication holds for every such price vector and deployment count. It does not cover a new policy or a controller refitted after the cost sample. A one-off resource ledger can use $K_k=\tau_k+\sum_jp_jW_{kj}$ with nonnegative prices on fitting, support construction, independent qualification and installation. When historical construction and current qualification were executed on different machines, their clocks remain distinct observations; pricing their individual resource components is not a claim to have measured a joint cold-start service. The deterministic support alone supplies no signed gain and cannot authorize purchase without the covered cost comparison.
